% SHARD-ID: AQM-38-LAGRANGIAN-STABILIZER-DESCENT
% SHARD-TITLE: Lagrangian Stabilizer Descent
% SHARD-SUMMARY: Defines the descent defect of a global stabilizer against a finite cover.
% SHARD-SUMMARY: Proves that Lagrangian labels give full stabilizer states, while zero descent defect is the extra local-to-global condition.
% SHARD-SUMMARY: Proves that disjoint products of local Lagrangians have zero descent defect.
% SHARD-KEYWORDS: Lagrangian, stabilizer, descent defect, local-to-global, cosheaf

\section{Lagrangian Stabilizer Descent}
\label{sec:lagrangian-stabilizer-descent}

\subsection{Status and Source Anchors}

This shard is a checked derivation from the already registered finite Weyl
and local-net layers.  AQM-08 and AQM-22 source and derive the prime-field
Weyl commutator and the criterion that stabilizer labels commute exactly when
their symplectic pairing vanishes.  AQM-26 proves that the Weyl label
assignment is the cosheaf-like additive layer and that the observable
assignment is a local net.  AQM-27 proves that quantum states do not form a
sheaf because local marginals do not determine global correlations.

The new convention is \texttt{CONVENTIONS.md} (ac).  We work here in the
finite vector-space case over an odd prime field \(k=\mathbb F_p\), with the
explicit counterexample over \(k=\mathbb F_3\).  Mixed residue fields and
finite Frobenius modules are not imported into the dimension statements in
this shard.

\subsection{Lamport Dependency Skeleton}

\[
\begin{array}{rcl}
D1&:&E=k^n_X\oplus k^n_Z,\quad
      \omega((x,z),(x',z'))=z\cdot x'-x\cdot z'.\\
D2&:&W(e)=T_xR_z\text{ on }\mathcal H=\ell^2(k^n),\quad
      W(e)W(f)=c(e,f)W(e+f).\\
D3&:&L\leq E\text{ is isotropic if }\omega(L,L)=0,
      \text{ and Lagrangian if }L=L^\perp.\\
D4&:&\alpha:L\to U(1)\text{ is a phase if }
      \alpha(e+f)=\alpha(e)\alpha(f)c(e,f)^{-1}.\\
L1&:&\{W(e):e\in E\}\text{ is a basis of }\operatorname{End}(\mathcal H).\\
L2&:&\text{A phased isotropic }L\text{ gives a commuting stabilizer group}.\\
T1&:&\text{A phased Lagrangian }L\text{ gives a rank-one full stabilizer state}.\\
D5&:&\operatorname{Loc}_U(L;\mathcal U)=
      \sum_i(L\cap E(U_i)),\quad
      D_U(L;\mathcal U)=L/\operatorname{Loc}_U(L;\mathcal U).\\
T2&:&D_U(L;\mathcal U)=0
      \Longleftrightarrow
      L\text{ is generated by stabilizers supported on }\mathcal U.\\
T3&:&\text{Disjoint products of local Lagrangians have zero defect.}
\end{array}
\]

\subsection{Finite Weyl Stabilizers D1--L2}

Let \(k=\mathbb F_p\) with \(p\) odd.  Put
\[
  E=k^n_X\oplus k^n_Z,\qquad
  \mathcal H=\ell^2(k^n),
\]
and write \(e=(x,z)\).  With a nontrivial additive character
\(\chi:k\to U(1)\), define
\[
  T_x|s\rangle=|s+x\rangle,\qquad
  R_z|s\rangle=\chi(z\cdot s)|s\rangle,\qquad
  W(x,z)=T_xR_z .
\]
As in AQM-22,
\[
  W(e)W(f)=\chi(z\cdot x')W(e+f),
  \qquad
  W(e)W(f)=\chi(\omega(e,f))W(f)W(e).
\]
Thus the multiplier in D2 is \(c(e,f)=\chi(z\cdot x')\).

\paragraph{Lemma \(L1\).}
The operators \(\{W(e):e\in E\}\) form a basis of
\(\operatorname{End}(\mathcal H)\).

\emph{Proof.}
There are \(|E|=p^{2n}\) such operators, and
\(\dim_\mathbb C\operatorname{End}(\mathcal H)=p^{2n}\).  It is enough to
prove linear independence.  For \(e=(x,z)\) and \(f=(x',z')\), the trace of
\(W(e)^*W(f)\) is zero if \(x\ne x'\), because the resulting translation has
no fixed basis vector.  If \(x=x'\) but \(z\ne z'\), the translation cancels
and the trace is
\[
  \sum_{s\in k^n}\chi((z'-z)\cdot s)=0
\]
by nontriviality of the character on the nonzero linear functional
\((z'-z)\cdot s\).  If \(e=f\), the trace is \(p^n\).  The Weyl operators are
therefore orthogonal and nonzero for the Hilbert-Schmidt pairing, hence
linearly independent.

\paragraph{Lemma \(L2\).}
Let \(L\leq E\) be isotropic and let \(\alpha:L\to U(1)\) satisfy D4.  Then
\[
  S_e=\alpha(e)^{-1}W(e),\qquad e\in L,
\]
is an ordinary unitary representation of the additive group \(L\).  In
particular the \(S_e\) commute.

\emph{Proof.}
Using D4 and the Weyl product law,
\[
\begin{aligned}
  S_eS_f
  &=\alpha(e)^{-1}\alpha(f)^{-1}c(e,f)W(e+f)\\
  &=\alpha(e+f)^{-1}W(e+f)\\
  &=S_{e+f}.
\end{aligned}
\]
The operators \(W(e)\) are unitary, and the phases have unit modulus, so each
\(S_e\) is unitary.  Since \(L\) is an Abelian additive group,
\(S_eS_f=S_{e+f}=S_{f+e}=S_fS_e\).

\subsection{Full Stabilizers Are Lagrangian T1}

\paragraph{Theorem \(T1\).}
If \(L\leq E\) is Lagrangian and phased by \(\alpha\), then the common
\(+1\)-eigenspace of the stabilizer group \(\{S_e:e\in L\}\) is one
dimensional.  Equivalently, a phased Lagrangian is a full stabilizer state at
the label level.

\emph{Proof.}
Because \(E\) is a \(2n\)-dimensional nondegenerate symplectic vector space,
\(\dim L+\dim L^\perp=2n\).  Thus \(L=L^\perp\) implies \(\dim L=n\), hence
\(|L|=p^n\).  The group average
\[
  P_L=\frac1{|L|}\sum_{e\in L}S_e
\]
is the orthogonal projection onto the common \(+1\)-eigenspace.  Indeed,
\[
  P_L^2=\frac1{|L|^2}\sum_{e,f\in L}S_{e+f}
  =
  \frac1{|L|}\sum_{g\in L}S_g=P_L,
\]
because every \(g\in L\) has exactly \(|L|\) presentations \(g=e+f\).
Also \(P_L^*=|L|^{-1}\sum_eS_{-e}=P_L\), and
\(S_fP_L=P_L=P_LS_f\) by the substitution \(e\mapsto f+e\).  Conversely, if
\(S_fv=v\) for all \(f\in L\), then \(P_Lv=v\).

Its trace is
\[
  \operatorname{Tr}(P_L)
  =
  \frac1{|L|}\sum_{e\in L}\alpha(e)^{-1}\operatorname{Tr}(W(e)).
\]
By Lemma \(L1\)'s trace calculation, \(\operatorname{Tr}(W(e))=0\) for
\(e\ne0\) and \(\operatorname{Tr}(W(0))=p^n\).  Therefore
\[
  \operatorname{Tr}(P_L)=\frac{p^n}{p^n}=1.
\]
An orthogonal projection has trace equal to its rank, so the stabilizer space
is one dimensional.

The same calculation also shows maximality among Weyl observables.  If an
operator \(A=\sum_{a\in E}\lambda_aW(a)\) commutes with every \(S_e\), then
each nonzero coefficient can occur only for labels \(a\in L^\perp=L\), because
the commutator factor is \(\chi(\omega(a,e))\).  Thus the centralizer of the
stabilizer group inside the Weyl basis is spanned by the same labels.

\subsection{Descent Defect D5--T2}

Now let \(U\) be a finite region or finite open in the residue-qudit net, and
let \(\mathcal U=\{U_i\}\) be a finite cover.  The label cosheaf of AQM-26
includes each \(E(U_i)\) into \(E(U)\) by extension by zero.  For a global
stabilizer label space \(L\leq E(U)\), define
\[
  \operatorname{Loc}_U(L;\mathcal U)
  =
  \sum_i\bigl(L\cap E(U_i)\bigr)
  \leq L
\]
and
\[
  D_U(L;\mathcal U)
  =
  L/\operatorname{Loc}_U(L;\mathcal U).
\]
This quotient is the stabilizer descent defect.

\paragraph{Theorem \(T2\).}
For a fixed global phased stabilizer \(L\leq E(U)\), the following are
equivalent:
\[
\begin{array}{ll}
(1)&D_U(L;\mathcal U)=0,\\
(2)&L=\sum_i(L\cap E(U_i)),\\
(3)&\{S_e:e\in L\}\text{ is generated by the stabilizers supported in the }
    U_i.
\end{array}
\]

\emph{Proof.}
The equivalence of (1) and (2) is exactly the definition of the quotient
defect.  If (2) holds, write any \(e\in L\) as \(e=e_1+\cdots+e_m\) with each
\(e_j\in L\cap E(U_{i_j})\).  Lemma \(L2\) gives
\[
  S_e=S_{e_1+\cdots+e_m}=S_{e_1}\cdots S_{e_m},
\]
so the global stabilizer is generated by local stabilizers.  Conversely, if
the stabilizer group is generated by supported stabilizers, then passing to
projective labels gives every \(e\in L\) as a sum of elements from
\(L\cap E(U_i)\), so (2) holds.

Therefore vanishing defect is the label-level local-to-global condition.  It
is the generation half of a cosheaf condition for a stabilizer subobject of
the label cosheaf.  Phase gluing is additional character compatibility on
overlaps; in the disjoint examples below there is no overlap condition.

\subsection{The Product Case T3}

\paragraph{Theorem \(T3\).}
Let \(U=U_1\sqcup U_2\), so
\[
  E(U)=E(U_1)\oplus E(U_2)
\]
as an orthogonal symplectic direct sum.  If \(L_i\leq E(U_i)\) is Lagrangian
for \(i=1,2\), then
\[
  L=L_1\oplus L_2\leq E(U)
\]
is Lagrangian and \(D_U(L;\{U_1,U_2\})=0\).  With product phases, the
stabilizer projection is \(P_{L_1}\otimes P_{L_2}\).

\emph{Proof.}
Orthogonality of the direct sum gives
\[
  L^\perp=L_1^\perp\oplus L_2^\perp=L_1\oplus L_2=L,
\]
so \(L\) is Lagrangian.  Also
\[
  L\cap E(U_1)=L_1,\qquad L\cap E(U_2)=L_2,
\]
hence \(\operatorname{Loc}_U(L;\{U_1,U_2\})=L\) and the defect vanishes.  The
group average over \(L_1\oplus L_2\) factors as the product of the two group
averages, which is \(P_{L_1}\otimes P_{L_2}\).

\subsection{Conclusion}

A Lagrangian subspace is the label-level form of a full stabilizer state.  A
vanishing stabilizer descent defect is the separate condition that this full
stabilizer is generated from local pieces of a chosen cover.  In cosheaf
language, \(E\) is the additive label cosheaf; a stabilizer choice has the
local-to-global generation property only when its labels have zero descent
defect.  The next shard gives the explicit two-qutrit counterexample showing
that the Lagrangian condition alone does not force this.
